\chapter{Government Bonds}\label{ch:m2:government-bonds}

Two traders on a chat agree to trade USD~10 million of a ten-year note at a
yield of 4.200\%. When the tickets come back they differ by USD~47\,351. Both
priced the same note at the same yield on the same settlement date, and both
are right: one booked the price quoted on the screen, 100-12+, and the other
the amount that changes hands, which adds the \omterm{def:m2:government-bonds:coupon}{coupon} interest earned since the
last payment date, forty-one days of a 4.25\% \omterm{def:m2:government-bonds:coupon}{coupon}. The note is the simplest
instrument in this book: a fixed schedule of payments from a government that
borrows in its own currency. Everything about it, its quoted price, its
yield, its sensitivity to rates, depends on conventions that are fixed by
regulation and market practice and that any system touching it must get
exactly right. This chapter sets them out and derives the risk measures on
which the rest of the book rests.

\section{Cash flows, day counts and accrued interest}

\begin{definition}[Coupon bond]\label{def:m2:government-bonds:coupon}
A fixed-coupon bond pays, per 100 of face value, a \emph{coupon}\index{coupon}
of $c/f$ on each of $f$ dates a year and repays its face value, called
\emph{par}\index{par}, at maturity. The coupon dates are generated backwards
from the maturity date in steps of $12/f$ months; if the maturity is the last
day of a month, so is every coupon date.
\end{definition}

\begin{definition}[Day-count convention]\label{def:m2:government-bonds:daycount}
A \emph{day-count convention}\index{day-count convention} turns two dates into
a fraction of a year. \emph{Actual/360} divides the actual number of days by
360; \emph{actual/365 fixed} by 365; \emph{30/360} counts every month as 30
days; \emph{actual/actual} (ICMA) counts the actual days elapsed in a \omterm{def:m2:government-bonds:coupon}{coupon}
period and divides by the actual days in that period, so every full period
earns exactly $c/f$.
\end{definition}

\begin{definition}[Accrued interest, clean and dirty price]\label{def:m2:government-bonds:accrued}
Between two \omterm{def:m2:government-bonds:coupon}{coupon} dates the seller of a bond has earned part of the next
\omterm{def:m2:government-bonds:coupon}{coupon}, and the buyer pays it to her: this is \emph{accrued
interest}\index{accrued interest}, $A = \frac{c}{f}\,(1 - w)$ under
actual/actual, where $w$ is the fraction of the current period still to run.
The \emph{dirty price}\index{dirty price} (or full, or invoice price) is what
changes hands; the \emph{clean price}\index{clean price} quoted on screens is
the dirty price less accrued interest.
\end{definition}

The market quotes \omterm{def:m2:government-bonds:accrued}{clean prices} because they do not jump on \omterm{def:m2:government-bonds:coupon}{coupon} dates
(\cref{fig:m2:government-bonds:cleandirty}): a \omterm{def:m2:government-bonds:accrued}{clean price} that moves means
that the market moved. In the dollar market \omterm{def:m2:government-bonds:accrued}{clean prices} of notes and bonds
are quoted in thirty-seconds of a point, so 100-12 is $100 + 12/32 =
100.375$, and a trailing plus adds a sixty-fourth: 100-12+ is 100.390625.

\begin{example}[Forty-one days]\label{ex:m2:government-bonds:accrued}
A note paying 4.25\% semiannually on 15 February and 15 August settles on 25
September 2026. The current period runs 184 days, of which 41 have elapsed:
$A = 2.125 \times 41/184 = 0.473505$ per 100, USD~4\,735 per million, USD~47\,351
on the trade of the opening paragraph.
\end{example}

\begin{omfigure}
\begin{tikzpicture}[font=\small]
\draw[thick,-{Stealth[length=2.5mm]}] (0,0) -- (13.2,0) node[right]{time};
\foreach \x/\t/\s in {1.0/previous coupon/15 Aug 2026,6.0/next coupon/15 Feb 2027,11.0/following coupon/15 Aug 2027}{
  \draw[thick] (\x,0.15) -- (\x,-0.15);
  \node[below=4pt,font=\footnotesize,align=center] at (\x,-0.15) {\t\\\s};}
\draw[omThm,very thick] (2.25,0.25) -- (2.25,-0.25);
\draw[omThm,densely dotted] (2.25,0.25) -- (2.25,1.1);
\node[anchor=south,font=\footnotesize\bfseries,text=omThm] at (2.25,1.1) {settlement, 25 Sep};
\draw[decorate,decoration={brace,amplitude=4pt},omProp,thick] (1.0,0.3) -- (2.25,0.3) node[midway,above=5pt,font=\footnotesize,text=omProp] {$1 - w$};
\draw[decorate,decoration={brace,amplitude=4pt},omDef,thick] (2.25,0.3) -- (6.0,0.3) node[midway,above=5pt,font=\footnotesize,text=omDef] {$w = 143/184$};
\foreach \x in {6.0,11.0}{\draw[omDef,-{Stealth[length=2mm]},thick] (\x,0.6) -- (\x,1.4);}
\node[font=\footnotesize,text=omDef,align=center] at (8.5,1.35) {coupons of $c/f$, discounted\\over $w$, $1 + w$, \dots\ periods};
\end{tikzpicture}

\omcaption{A settlement between \omterm{def:m2:government-bonds:coupon}{coupon} dates. The seller is paid the
fraction $1 - w$ of the coming \omterm{def:m2:government-bonds:coupon}{coupon} as \omterm{def:m2:government-bonds:accrued}{accrued interest}; the buyer receives
the whole \omterm{def:m2:government-bonds:coupon}{coupon}. The street convention discounts each later flow over $w$,
$1 + w$, $2 + w$, \dots\ periods.}\label{fig:m2:government-bonds:timeline}
\end{omfigure}

\begin{dated}{2026-09}{Conventions of three sovereign markets}\label{dat:m2:government-bonds:conventions}
\textbf{US Treasury notes and bonds}: semiannual \omterm{def:m2:government-bonds:coupon}{coupons} on dates set by the
maturity; each regular payment is exactly half the annual \omterm{def:m2:government-bonds:coupon}{coupon}; \omterm{def:m2:government-bonds:accrued}{accrued
interest} on the actual days of the half-year; prices in 32nds of a point.
\textbf{UK gilts}: semiannual \omterm{def:m2:government-bonds:coupon}{coupons}; actual/actual \omterm{def:m2:government-bonds:accrued}{accrued interest} since
November 1998; ex-dividend seven business days before each payment, after
which the buyer does not receive the coming \omterm{def:m2:government-bonds:coupon}{coupon}. \textbf{German Bunds}:
issued at 7, 10, 15 or 30 years with a fixed \emph{annual} \omterm{def:m2:government-bonds:coupon}{coupon}.
\end{dated}

\section{Price and yield}

\begin{definition}[Yield to maturity]\label{def:m2:government-bonds:ytm}
The \emph{yield to maturity}\index{yield to maturity} $y$ of a bond with $n$
remaining \omterm{def:m2:government-bonds:coupon}{coupons}, settled at fraction $w$ of a period before the next one,
is the rate, compounded $f$ times a year, that equates the \omterm{def:m2:government-bonds:accrued}{dirty price} with
its discounted flows:
\[
P_{\text{dirty}} \;=\; \sum_{k=0}^{n-1} \frac{c/f}{(1 + y/f)^{k+w}} \;+\; \frac{100}{(1 + y/f)^{n-1+w}} .
\]
\end{definition}

\begin{proposition}[Price and yield]\label{prop:m2:government-bonds:priceyield}
The \omterm{def:m2:government-bonds:accrued}{dirty price} is a strictly decreasing, strictly convex function of $y$. On
a \omterm{def:m2:government-bonds:coupon}{coupon} date ($w = 1$ for the period just starting) the price is 100 if and
only if $y = c/100$.
\end{proposition}
\begin{proof}
Each term $a_k (1 + y/f)^{-t_k}$ with $a_k, t_k > 0$ is strictly decreasing
and strictly convex in $y$, so is their sum. With $w = 1$ and $v = 1/(1 +
y/f)$ the price is $\frac cf \sum_{k=1}^{n} v^k + 100 v^n$; with $y/f =
c/(100 f)$ the annuity equals $(1 - v^n)/(y/f)$ and the price is $100(1 - v^n)
+ 100 v^n = 100$. Uniqueness follows from monotonicity.
\end{proof}

\begin{remark}[Two conventions for the fraction]\label{rem:m2:government-bonds:conventions}
The formula above, which compounds over the fraction $w$, is the market's.
The US Treasury's own regulation, used to set auction prices, discounts the
fractional period at simple interest, $P_{\text{dirty}}(1 + w\,y/f) =$ value at
the next \omterm{def:m2:government-bonds:coupon}{coupon} date. The two agree on \omterm{def:m2:government-bonds:coupon}{coupon} dates and differ by a few
thousandths of a point between them. A system that reproduces the
regulation's worked examples and the market's screens must implement both
(\cref{bld:m2:government-bonds:bond}).
\end{remark}

\begin{dated}{2026-09}{US marketable debt}\label{dat:m2:government-bonds:debt}
Outstanding on 31 August 2026: USD~31.8 trillion of marketable Treasury
securities, of which bills 7.2 trillion, notes 16.2 trillion, bonds 5.5
trillion, inflation-protected securities 2.2 trillion and floating-rate notes
0.7 trillion.
\end{dated}

\section{Duration, DV01 and convexity}

\begin{definition}[Duration, DV01, convexity]\label{def:m2:government-bonds:duration}
With flows $a_k$ at times $t_k$ (in periods) and $v = 1/(1 + y/f)$, the
\emph{Macaulay duration}\index{Macaulay duration} is the present-value
weighted average time of the flows, $D = \frac{1}{f P}\sum_k t_k a_k v^{t_k}$,
in years. The \emph{modified duration}\index{modified duration} is
$D_{\mathrm{mod}} = -\frac1P \frac{dP}{dy} = D/(1 + y/f)$. The
\emph{DV01}\index{DV01} is the fall in price for a one-basis-point rise in
yield, $P\,D_{\mathrm{mod}} \times 10^{-4}$ per 100, usually quoted per million
of face. The \emph{convexity}\index{convexity} is $\mathcal C = \frac1P
\frac{d^2P}{dy^2}$.
\end{definition}

\begin{proposition}[Second-order price change]\label{prop:m2:government-bonds:taylor}
For a parallel change $\Delta y$ of the yield,
\[
\frac{\Delta P}{P} \;=\; -D_{\mathrm{mod}}\,\Delta y \;+\; \tfrac12\,\mathcal C\,(\Delta y)^2 \;+\; O(\Delta y^3),
\qquad
\mathcal C \;=\; \frac{1}{f^2 P (1+y/f)^2}\sum_k t_k(t_k+1)\,a_k v^{t_k} .
\]
\end{proposition}
\begin{proof}
Taylor's formula in $y$, with $\frac{d}{dy} v^{t} = -\frac{t}{f}\,v^{t+1}$ and
$\frac{d^2}{dy^2} v^t = \frac{t(t+1)}{f^2}\, v^{t+2}$.
\end{proof}

\begin{example}[The ten-year note]\label{ex:m2:government-bonds:risk}
At 4.200\% the note of \cref{ex:m2:government-bonds:accrued} has a \omterm{def:m2:government-bonds:accrued}{clean price}
of 100.397405 (100-12+), a \omterm{def:m2:government-bonds:accrued}{dirty price} of 100.870911, a \omterm{def:m2:government-bonds:duration}{Macaulay duration} of
8.142 years, a \omterm{def:m2:government-bonds:duration}{modified duration} of 7.975, a \omterm{def:m2:government-bonds:duration}{DV01} of USD~804 per million and a
\omterm{def:m2:government-bonds:duration}{convexity} of 75.8. A 100-basis-point rise costs 7.675 points; duration alone
predicts 8.044, duration and \omterm{def:m2:government-bonds:duration}{convexity} together 7.662
(\cref{fig:m2:government-bonds:priceyield}).
\end{example}

\begin{omfigure}
\begin{tikzpicture}
\begin{axis}[width=11cm,height=5.2cm,
  xlabel={yield (\%)},ylabel={dirty price},
  tick label style={font=\small},label style={font=\small},
  xmin=1.7,xmax=6.7,ymin=80,ymax=126,grid=major,grid style={gray!25},
  legend columns=3,legend style={font=\footnotesize,at={(0.5,-0.3)},anchor=north,draw=none,fill=none,/tikz/every even column/.append style={column sep=8pt}},
  table/col sep=comma]
\addplot[omDef,very thick] table[x=yield,y=price]{figdata/markets-2/03-government-bonds/priceyield.csv};
\addplot[omThm,thick,dashed] table[x=yield,y=duration]{figdata/markets-2/03-government-bonds/priceyield.csv};
\addplot[omProp,thick,dotted] table[x=yield,y=convexity]{figdata/markets-2/03-government-bonds/priceyield.csv};
\legend{price,duration (tangent),duration and convexity}
\end{axis}
\end{tikzpicture}

\omcaption{Price against yield for the ten-year note of
\cref{ex:m2:government-bonds:risk}, two and a half percentage points either
side of 4.20\%. The tangent (duration) lies below the curve on both sides: a
bond gains more when yields fall than it loses when they rise by the same
amount, here by 2.2 to 2.6 points at the edges. The second-order approximation
stays close to the curve over the whole range. Data: the chapter's tutorial.}\label{fig:m2:government-bonds:priceyield}
\end{omfigure}

\begin{method}[Hedging one bond with another]\label{met:m2:government-bonds:hedge}
To hedge a face $N_1$ of bond 1 against parallel moves with bond 2, sell $N_2
= N_1\,\mathrm{DV01}_1/\mathrm{DV01}_2$ of face. The hedge leaves: the
difference in \omterm{def:m2:government-bonds:duration}{convexity}, the risk that the two yields do not move together
(curve risk), and the drift of both \omterm{def:m2:government-bonds:duration}{DV01s} as time passes and yields move,
which calls for rebalancing.
\end{method}

\section{Zero-coupon rates, par yields and strips}

\begin{definition}[Zero-coupon rate, par yield, STRIPS]\label{def:m2:government-bonds:zero}
The \emph{zero-coupon rate}\index{zero-coupon rate} $z(T)$ for maturity $T$
is the yield of a single payment at $T$: $P(0,T) = (1 + z/f)^{-fT}$. The
\emph{par yield}\index{par yield} for $T$ is the \omterm{def:m2:government-bonds:coupon}{coupon} that would price a
bond maturing at $T$ at 100. \emph{STRIPS}\index{STRIPS} are the US
Treasury's zero-coupon securities, created by separating an eligible note or
bond into its individual \omterm{def:m2:government-bonds:coupon}{coupon} and principal payments, each of which then
trades on its own; a complete set can be reassembled into the original
security.
\end{definition}

\begin{proposition}[Par yields from discount factors, and back]\label{prop:m2:government-bonds:par}
With discount factors $P_k = P(0, k/f)$,
\[
\text{par}(T_n) \;=\; f\,\frac{1 - P_n}{\sum_{k=1}^n P_k},
\qquad
P_n \;=\; \frac{1 - \frac{\text{par}_n}{f}\sum_{k<n} P_k}{1 + \frac{\text{par}_n}{f}} .
\]
The second formula \emph{bootstraps} the discount factors one maturity at a
time from \omterm{def:m2:government-bonds:zero}{par yields}.
\end{proposition}
\begin{proof}
A bond with \omterm{def:m2:government-bonds:coupon}{coupon} $c$ at \omterm{def:m2:government-bonds:coupon}{par} satisfies $100 = \frac cf \sum_{k \le n} P_k +
100 P_n$; solve for $c/100$, or, given $c/100 = \text{par}_n$ and $P_1, \dots,
P_{n-1}$, for $P_n$.
\end{proof}

\begin{omfigure}
\begin{tikzpicture}
\begin{axis}[width=11cm,height=4.8cm,
  xlabel={maturity (years)},ylabel={rate (\%)},
  tick label style={font=\small},label style={font=\small},
  xmin=0,xmax=10.5,ymin=3.88,ymax=4.26,grid=major,grid style={gray!25},
  legend columns=2,legend style={font=\footnotesize,at={(0.5,-0.3)},anchor=north,draw=none,fill=none,/tikz/every even column/.append style={column sep=8pt}},
  table/col sep=comma]
\addplot[omDef,very thick,mark=*,mark size=1.2pt] table[x=years,y=par]{figdata/markets-2/03-government-bonds/zero.csv};
\addplot[omThm,very thick,dashed] table[x=years,y=zero]{figdata/markets-2/03-government-bonds/zero.csv};
\legend{par yield,zero-coupon rate (bootstrapped)}
\end{axis}
\end{tikzpicture}

\omcaption{An illustrative \omterm{def:m2:government-bonds:coupon}{par} curve and the zero curve bootstrapped from it
(\cref{prop:m2:government-bonds:par}). Where the \omterm{def:m2:government-bonds:coupon}{par} curve rises, the zero
curve lies above it, because a \omterm{def:m2:government-bonds:coupon}{coupon} bond's yield averages the zero rates of
all its payments, most of them earlier and lower; where the \omterm{def:m2:government-bonds:coupon}{par} curve dips,
at one to two years, the zero curve dips below it. Data: the chapter's
tutorial.}\label{fig:m2:government-bonds:zero}
\end{omfigure}

\begin{omfigure}
\begin{tikzpicture}
\begin{axis}[width=11cm,height=4.8cm,
  xlabel={days after 15 August 2026},ylabel={price},
  tick label style={font=\small},label style={font=\small},
  xmin=0,xmax=364,ymin=100.2,ymax=102.8,grid=major,grid style={gray!25},
  legend columns=2,legend style={font=\footnotesize,at={(0.5,-0.3)},anchor=north,draw=none,fill=none,/tikz/every even column/.append style={column sep=8pt}},
  table/col sep=comma]
\addplot[omThm,very thick] table[x=day,y=dirty]{figdata/markets-2/03-government-bonds/cleandirty.csv};
\addplot[omDef,very thick] table[x=day,y=clean]{figdata/markets-2/03-government-bonds/cleandirty.csv};
\legend{dirty price,clean price}
\end{axis}
\end{tikzpicture}

\omcaption{The ten-year note over a year at a constant yield of 4.20\%, daily.
The \omterm{def:m2:government-bonds:accrued}{dirty price} climbs as the \omterm{def:m2:government-bonds:coupon}{coupon} accrues and drops by the \omterm{def:m2:government-bonds:coupon}{coupon} when it
is paid (15 February, day 184); the \omterm{def:m2:government-bonds:accrued}{clean price} moves only by its slow pull
towards \omterm{def:m2:government-bonds:coupon}{par}. Data: the chapter's tutorial.}\label{fig:m2:government-bonds:cleandirty}
\end{omfigure}

\section{Tutorial: a bond on a real calendar}\label{tut:m2:government-bonds:bond}

\begin{tutorial}
\textbf{Goal.} Price the ten-year note from its yield on 25 September 2026,
check its \omterm{def:m2:government-bonds:duration}{DV01} against a bumped reprice, and bootstrap a zero curve.
\textbf{End state:} the numbers of \cref{ex:m2:government-bonds:risk} and
\cref{fig:m2:government-bonds:priceyield,fig:m2:government-bonds:zero}.
\begin{tutsteps}
\item \textbf{The schedule and the \omterm{def:m2:government-bonds:accrued}{accrued interest}}, generated back from
maturity with the end-of-month rule.
\omcode{code/firm/bond/firm_bond.py}{51}{71}{Coupon dates, the fraction of the period to run, accrued interest.}\label{lst:m2:government-bonds:schedule}
You should see 0.473505 for the note on 25 September 2026.
\item \textbf{The price} in both conventions of \cref{rem:m2:government-bonds:conventions}.
\omcode{code/firm/bond/firm_bond.py}{73}{78}{Dirty price from yield: street and Treasury conventions.}\label{lst:m2:government-bonds:dirty}
Check the regulation's examples: 99.057893 and 99.730918.
\item \textbf{Risk}, in the systems language of the firm. The C++20 twin
computes duration, \omterm{def:m2:government-bonds:duration}{DV01} and \omterm{def:m2:government-bonds:duration}{convexity} in one pass over the flows.
\omcode{code/firm/bond/cpp/firm_bond.hpp}{94}{108}{Duration, DV01 and convexity (C++20).}\label{lst:m2:government-bonds:risk}
\omterm{def:m2:government-bonds:duration}{DV01} should equal half the price difference between yields 1 basis point
either side, to six significant figures.
\item \textbf{The zero curve.}
\omcode{code/firm/bond/firm_bond.py}{116}{122}{Bootstrapping discount factors from par yields.}\label{lst:m2:government-bonds:bootstrap}
\texttt{fig\_bond.py} writes the three charts of the chapter.
\end{tutsteps}
\textbf{What to change next.} Price the same note on 13 February 2027 and on
16 February 2027 and compare clean and \omterm{def:m2:government-bonds:accrued}{dirty prices} across the \omterm{def:m2:government-bonds:coupon}{coupon}; then
set the \omterm{def:m2:government-bonds:coupon}{coupon} to zero and check that the \omterm{def:m2:government-bonds:duration}{Macaulay duration} equals the
remaining life.
\end{tutorial}

\section{Build: the bond library}\label{bld:m2:government-bonds:bond}

\begin{build}
\bfield{Purpose} The first piece of the miniature firm's pricing library
(extended in One Quant Books 5 and 6): every rates component, the futures basis
(\cref{ch:m2:bond-futures}), the auction analyser
(\cref{ch:m2:the-treasury-market}), the spread measures of the credit part
(\cref{ch:m2:corporate-bonds}), prices bonds through it. Python for research,
C++20 for the trading path, Rust as the twin.
\bfield{Interface} \texttt{Bond(coupon, maturity, freq)} with
\texttt{coupon\_dates(settle)}, \texttt{accrued(settle)},
\texttt{dirty\_price(y, settle, treasury)}, \texttt{clean\_price(\ldots)},
\texttt{yield\_from\_clean(clean, settle)}, \texttt{risk(y, settle)};
\texttt{act\_360}, \texttt{act\_365f}, \texttt{thirty\_360};
\texttt{bootstrap\_par(pars)}, \texttt{par\_yield(dfs)}. The C++ header and the
Rust crate expose the same names.
\bfield{Rules} \omterm{def:m2:government-bonds:coupon}{Coupon} dates from maturity backwards, end-of-month preserved;
actual/actual \omterm{def:m2:government-bonds:accrued}{accrued interest}; the street convention by default and the
Treasury convention on request; yields decimal, compounded at the \omterm{def:m2:government-bonds:coupon}{coupon}
frequency.
\bfield{Acceptance tests} \texttt{code/firm/bond/tests/},
\texttt{cpp/firm\_bond\_test.cpp}, \texttt{rust/}: the regulation's worked
examples (99.057893; accrued 0.367403 and 99.730918); a leap-year month-end
schedule; price and yield round trip; \omterm{def:m2:government-bonds:duration}{DV01} and \omterm{def:m2:government-bonds:duration}{convexity} against bumps; a zero
bond's duration equals its life; bootstrapped factors reprice their \omterm{def:m2:government-bonds:zero}{par
yields}.
\bfield{Stretch} Short and long first \omterm{def:m2:government-bonds:coupon}{coupons} (the regulation's other
examples); ex-dividend periods as in gilts; the 30/360 European variant; a
business-day calendar for payment dates.
\end{build}

\begin{omsources}
\item 31 CFR Part 356, Appendix B, \emph{Formulas and Tables}, sections I and
II. US Treasury, \emph{STRIPS} (TreasuryDirect); \emph{Monthly Statement of
the Public Debt}, August 2026.
\item UK Debt Management Office, \emph{Formulae for Calculating Gilt Prices
from Yields}. Deutsche Finanzagentur, \emph{Federal Bonds}.
\item P.\ Ritchken, \emph{Fixed Income}, chapter 2, ``Treasury Securities''.
\item F.\,J.\ Fabozzi (ed.), \emph{The Handbook of Fixed Income Securities},
McGraw-Hill, for the conventions of other markets.
\end{omsources}

\section{Exercises}

\begin{exercise}[$\star$]\label{exo:m2:government-bonds:1}
A 4.25\% semiannual note paid a \omterm{def:m2:government-bonds:coupon}{coupon} on 15 August. Give its \omterm{def:m2:government-bonds:accrued}{accrued interest}
per 100 for settlement on 25 September, the period being 184 days long.
\end{exercise}

\begin{exercise}[$\star$]\label{exo:m2:government-bonds:2}
Convert 99-16+ to a decimal price, and 101.296875 to a quote in 32nds.
\end{exercise}

\begin{exercise}[$\star$]\label{exo:m2:government-bonds:3}
Give the \omterm{def:m2:government-bonds:accrued}{clean price}, on a \omterm{def:m2:government-bonds:coupon}{coupon} date, of a three-year 4\% semiannual bond
at a yield of 4\%, then at 5\%.
\end{exercise}

\begin{exercise}[$\star\star$]\label{exo:m2:government-bonds:4}
For a two-year 4\% semiannual bond at a yield of 4\% on a \omterm{def:m2:government-bonds:coupon}{coupon} date, give
the \omterm{def:m2:government-bonds:duration}{Macaulay duration}, the \omterm{def:m2:government-bonds:duration}{modified duration} and the \omterm{def:m2:government-bonds:duration}{DV01} per million.
\end{exercise}

\begin{exercise}[$\star\star$]\label{exo:m2:government-bonds:5}
For the ten-year note of \cref{ex:m2:government-bonds:risk}, estimate the
price change for a 100-basis-point \emph{fall} in yield with duration alone
and with \omterm{def:m2:government-bonds:duration}{convexity}, and compare with a full reprice.
\end{exercise}

\begin{exercise}[$\star\star$]\label{exo:m2:government-bonds:6}
A 5\% semiannual bond pays on 15 February and 15 August. Compute its \omterm{def:m2:government-bonds:accrued}{accrued
interest} for settlement on 31 March 2026 under actual/actual and under 30/360.
Which market's investor would care about the difference?
\end{exercise}

\begin{exercise}[$\star\star\star$]\label{exo:m2:government-bonds:7}
\emph{Coding.} With \texttt{bootstrap\_par}, bootstrap the \omterm{def:m2:government-bonds:coupon}{par} curve of
\cref{fig:m2:government-bonds:zero} and report the ten-year discount factor and
zero-coupon rate. Explain the sign of zero minus \omterm{def:m2:government-bonds:coupon}{par} at ten years.
\end{exercise}

\begin{exercise}[$\star\star\star$]\label{exo:m2:government-bonds:8}
\emph{Find the flaw.} ``A new thirty-year bond has a duration of thirty
years, so a one-basis-point rise in yields costs 0.30\% of its price.''
Correct it for a 4.5\% semiannual bond issued at \omterm{def:m2:government-bonds:coupon}{par}.
\end{exercise}

\section{Problem: The Two-Bond Desk}

\begin{problem}[{Weekend problem --- a DV01-neutral curve position}]\label{pb:m2:government-bonds:1}
On 25 September 2026 a desk buys USD~100 million face of a five-year note (4\%
\omterm{def:m2:government-bonds:coupon}{coupon}, maturing 30 September 2031, yield 3.95\%) and hedges it with the ten-year
note of \cref{ex:m2:government-bonds:risk} (4.25\%, 15 August 2036, yield
4.20\%).

\medskip\noindent\textbf{Part I --- Prices.}
\begin{enumerate}
\item Give the five-year note's previous and next \omterm{def:m2:government-bonds:coupon}{coupon} dates. Why is its
next \omterm{def:m2:government-bonds:coupon}{coupon} so close?
\item Give its \omterm{def:m2:government-bonds:accrued}{accrued interest} and its \omterm{def:m2:government-bonds:accrued}{clean price}, in decimal and in 32nds.
\item Give its dirty value for USD~100 million face.
\item Give both notes' \omterm{def:m2:government-bonds:duration}{DV01} per million.
\item Give both notes' \omterm{def:m2:government-bonds:duration}{convexity}.
\end{enumerate}

\medskip\noindent\textbf{Part II --- The hedge.}
\begin{enumerate}[resume]
\item What face of the ten-year note must the desk sell to be DV01-neutral?
\item What is the net cash of the position (long dirty value less short dirty
value)?
\item Reprice both notes with both yields 1 basis point higher. What is the
P\&L?
\item Reprice with both yields 25 basis points higher, then 25 lower. What is
the P\&L each time?
\item Why is it negative both ways?
\end{enumerate}

\medskip\noindent\textbf{Part III --- The curve.}
\begin{enumerate}[resume]
\item The five-year yield rises 10 basis points, the ten-year does not move.
What is the P\&L?
\item The ten-year yield rises 10 basis points, the five-year does not move.
What is the P\&L?
\item Which curve move does the position profit from? Name it.
\item Is the position's duration zero? Its \omterm{def:m2:government-bonds:duration}{DV01}?
\item Four weeks later neither yield has moved. Is the position still
DV01-neutral? What changed?
\end{enumerate}

\medskip\noindent\textbf{Part IV --- Judgement.}
\begin{enumerate}[resume]
\item Why hedge \omterm{def:m2:government-bonds:duration}{DV01} rather than face value, or market value?
\item Name two risks the hedge leaves.
\item Why would the desk quote both notes in 32nds but compute in decimals?
\item State the \emph{named result}: the hedge ratio and the P\&L of the
hedged position for a 25-basis-point parallel move either way.
\item In one sentence: what does a DV01-neutral position bet on?
\end{enumerate}
\end{problem}

\section{Interview questions}

\begin{interviewq}[$\star$ \iqroles{trader, bank}]\label{iq:m2:government-bonds:1}
What is the difference between a bond's clean and \omterm{def:m2:government-bonds:accrued}{dirty price}, and why does
the market quote the clean one?
\end{interviewq}

\begin{interviewq}[$\star$ \iqroles{trader, researcher}]\label{iq:m2:government-bonds:2}
Why is the price--yield relation convex? Is \omterm{def:m2:government-bonds:duration}{convexity} good for the holder?
\end{interviewq}

\begin{interviewq}[$\star\star$ \iqroles{trader, researcher, bank}]\label{iq:m2:government-bonds:3}
A ten-year zero-coupon bond and a ten-year 5\% \omterm{def:m2:government-bonds:coupon}{coupon} bond: which has the
larger duration? The larger \omterm{def:m2:government-bonds:duration}{DV01} per 100 of face?
\end{interviewq}

\begin{interviewq}[$\star\star$ \iqroles{trader}]\label{iq:m2:government-bonds:4}
You are long the five-year and want to hedge with the ten-year. How much do
you sell, and what risks remain?
\end{interviewq}

\begin{interviewq}[$\star\star$ \iqroles{developer}]\label{iq:m2:government-bonds:5}
You must implement \omterm{def:m2:government-bonds:accrued}{accrued interest} for government bonds in three markets.
What will your tests cover?
\end{interviewq}

\begin{interviewq}[$\star\star\star$ \iqroles{researcher, trader}]\label{iq:m2:government-bonds:6}
What is the \omterm{def:m2:government-bonds:duration}{modified duration} of a perpetual bond paying a fixed \omterm{def:m2:government-bonds:coupon}{coupon} once
a year, at a yield of 4\%? And its \omterm{def:m2:government-bonds:duration}{Macaulay duration}?
\end{interviewq}
